\section{The no-go theorems and their hypotheses}
\label{sec:theorems}

We state each theorem in the form in which it is proved, and list its hypotheses with the labels
H0 to H8 that the rest of the paper uses (Table~\ref{tab:hypotheses}). A \emph{symmetry} of a
quantum theory with a scattering matrix is a linear operator on the Hilbert space of the theory
that commutes with the scattering matrix; a \emph{generator} of a one-parameter group of such
operators is the operator whose exponential gives the group. The \emph{Poincar\'e group} of a
spacetime is the group of its translations and Lorentz transformations; its generators are the
momenta and the generators of rotations and boosts.

\emph{Coleman and Mandula~\cite{ColemanMandula1967}} (in the form of Ref.~\cite{WeinbergQFT3},
Chap.~24). Let the symmetries of the scattering matrix form a connected Lie group of unitary operators
on a Hilbert space with positive scalar product (H0). Assume: (H1) the group contains the Poincar\'e
group of a spacetime with one time direction; (H2) for every mass there are finitely many kinds of
particles below it; (H3) the amplitudes of elastic two-particle scattering are analytic functions of
the energy and the momentum transfer near the physical region, apart from the normal thresholds;
(H4) two particles scatter at almost every energy and angle; (H5) the generators act on states of one
particle as integral operators in momentum space whose kernels are distributions, map one-particle
states to one-particle states, and act on a state of several particles as the sum of their action
on each particle. Then the group is locally the direct product of the Poincar\'e group and a group
whose generators commute with all momenta, rotations and boosts and form the direct sum of a compact
semisimple Lie algebra and an abelian one (for massless particles the Poincar\'e group may grow into
the conformal group). Pelc and Horwitz proved the same conclusion for any number of space dimensions
with one time~\cite{PelcHorwitz1997}; the theorem therefore applies to the $7+1$ spacetime, while on
the $4+4$ spacetime, with four time-like directions, its hypothesis H1 fails from the outset.

\emph{McGlinn~\cite{McGlinn1964}.} Let a Lie algebra be the sum of the Poincar\'e algebra and a set of
generators that commute with all momenta. If these generators form a semisimple Lie algebra (H6),
they commute with the rotations and boosts as well.

\emph{O'Raifeartaigh~\cite{ORaifeartaigh1965}.} In a unitary representation of a Lie algebra of finite
dimension that contains the Poincar\'e algebra, if the square of the momentum has an isolated
eigenvalue on an irreducible representation (H7), every state of that representation has the same
mass: a finite-dimensional symmetry algebra cannot split the masses of a multiplet.

\emph{Haag, {\L}opusza\'nski and Sohnius~\cite{HLS1975}, Nahm~\cite{Nahm1978}.} If the symmetry
generators may anticommute (a Lie superalgebra: even generators obey commutators, odd generators
anticommutators), the hypotheses H0 to H5 admit exactly one extension of the conclusion of Coleman
and Mandula: odd generators $\mathcal Q_{\rm NG}$ (supercharges) whose anticommutators give the
momenta, with an internal ``R-symmetry'' that rotates the supercharges, and central or tensorial
charges. Their proof uses the positivity of the anticommutator $\{\mathcal Q_{\rm NG},
\mathcal Q_{\rm NG}^\dagger\}$ in a Hilbert space with positive scalar product (H8); in this
extension Lorentz and internal symmetry are linked, because the supercharges carry spin and internal
charge at once.

\emph{Weinberg and Witten~\cite{WeinbergWitten1980}.} Massless particles of helicity larger than $1/2$
in magnitude cannot carry a conserved current that transforms as a Lorentz vector, and those of
helicity larger than 1 cannot carry a conserved energy-momentum tensor that transforms as a tensor.
The proof uses H0, H1 and the fact that the little group of a massless momentum (the subgroup of the
Lorentz group that leaves it unchanged) acts on helicity through a compact rotation group.

\begin{table}[t]
\caption{The hypotheses of the no-go theorems, with the labels used in this paper.}
\label{tab:hypotheses}
\begin{ruledtabular}
\begin{tabular}{l>{\raggedright\arraybackslash}p{0.38\textwidth}}
label & hypothesis\\
\hline
H0 & symmetries are linear operators on a Hilbert space with positive scalar product; their generators form a Lie algebra (an associative product, hence the Jacobi identity)\\
H1 & the scattering matrix and the vacuum are invariant under the Poincar\'e group of a spacetime with one time\\
H2 & finitely many kinds of particles below every mass\\
H3 & analyticity of elastic two-particle amplitudes\\
H4 & scattering occurs at almost all energies and angles\\
H5 & generators map one particle to one particle and act additively on several\\
H6 & (McGlinn) the internal algebra is semisimple\\
H7 & (O'Raifeartaigh) the squared momentum has an isolated eigenvalue\\
H8 & (Haag, {\L}opusza\'nski, Sohnius) the anticommutator of a supercharge with its adjoint is positive\\
\end{tabular}
\end{ruledtabular}
\end{table}
